\scanpage{274}
\mspage{15}

\deleted{Ceci rend immédiate la}

\deleted{\uline{Proposition 7} Soit $s\in I$ un élément qui}

\uline{Exemples}

1. $I=\Delta_0=e$, $\underline{\mathrm{Ch}}_I(X)\simeq X$. Les chemins de longueur $0$ s'identifient aux objets \inserted{$A$} de $X$. L'unique élément $\mathrm{id}_e$ de $\mathrm{Fl}(I)$ \deleted{\unsure{dev.|dev}} donne $\mathrm{id}_A$ dans $X$.

2. $I=\mathbb{\Delta}_1^+$, $\underline{\mathrm{Ch}}_I(X)\simeq\underline{\mathrm{Fl}}(X)$. Les chemins de \deleted{longueur $1$} type $\mathbb{\Delta}_1^+$ sont \inserted{ceux de longueur $1$ et \uline{directs}}, i.e. donnés par une flèche

\diagram{\begin{tikzcd}[ampersand replacement=\&]
a_0\arrow[r] \& a_1
\end{tikzcd}}

\deleted{qui est \unsure{la|le}} \inserted{dont la} source est aussi la source $a_0$ du chemin. Par contre, \deleted{\unsure{p.}} les chemins de type $\mathbb{\Delta}_1^-$ sont donnés par des diagrammes

\diagram{\begin{tikzcd}[ampersand replacement=\&]
a_0 \& a_1\arrow[l]
\end{tikzcd}}

i.e par des flèches dont le but est considéré comme source du chemin. Bien sûr,

% slanted
\margin{\unsure{?}}

\deleted{\unsure{Ch}} les catégories $\underline{\mathrm{Ch}}_{\mathbb{\Delta}_{\unsure{1}}^{+}}(X)$ et $\underline{\mathrm{Ch}}_{\mathbb{\Delta}_{\unsure{1}}^{-}}(X)$ sont canoniquement isomorphes, mais par un isomorphisme qui \uline{échange source et but \unsure{du|des} \unsure{chemins|chemin}}. Nous nous garderons donc bien de confondre les deux chemins précédents.

% archive sheet number: 273
